% A cone generated by two vectors (convex), and a cone that is not convex.
%
%   figures/tikz/ra-cone.tex
%       -> media/figures/ra-cone.{png,svg}
%
% Lecture 13 (Real Analysis Review), Cones.
%
% GEOMETRY (1 unit = 1 cm), window [-2.3,2.3]^2.
%  Left: d_1 = (1.8,0.5), d_2 = (0.6,1.7). C = {t_1 d_1 + t_2 d_2 : t_1,t_2 >= 0}
%   is the wedge between the rays through d_1 and d_2; convex. The point drawn
%   is x = 0.5 d_1 + 0.5 d_2 = (1.2,1.1), and 1.5 x = (1.8,1.65) shows a positive
%   multiple staying in C.
%  Right: C = {x : x_1 x_2 >= 0} (first and third quadrants), a cone since
%   (t x_1)(t x_2) = t^2 x_1 x_2 >= 0. Not convex: p = (1.6,0.4) and q = (-0.4,-1.6)
%   are in C, the midpoint (0.6,-0.6) has x_1 x_2 = -0.36 < 0, so it is not.
%
% GREYSCALE: the wedge(s) are hatched; the failing midpoint carries a drawn X
% (two crossing rules) and the words "not in C".
\usetikzlibrary{patterns}
\definecolor{okabeblue}{HTML}{0072B2}
\definecolor{okabevermillion}{HTML}{D55E00}

\begin{tikzpicture}[
    >={Latex[length=2.0mm]},
    font=\small,
    pt/.style={fill=black, draw=black, circle, inner sep=0pt, minimum size=3.4pt},
    lab/.style={font=\scriptsize, inner sep=1.5pt},
    panelcap/.style={font=\small, align=center, text width=60mm, anchor=north},
    ax/.style={draw=black!40, line width=0.4pt, ->},
  ]
% ------------------------------------------------------------- left
\begin{scope}
  \begin{scope}
    \clip (-2.3,-2.3) rectangle (2.3,2.3);
    \path[pattern=north east lines, pattern color=okabeblue!55]
       (0,0) -- (9.0,2.5) -- (3.0,8.5) -- cycle;
    \draw[okabeblue, line width=1.2pt] (0,0) -- (9.0,2.5);
    \draw[okabeblue, line width=1.2pt] (0,0) -- (3.0,8.5);
  \end{scope}
  \draw[ax] (-2.3,0) -- (2.4,0);
  \draw[ax] (0,-2.3) -- (0,2.4);
  \draw[->, line width=1.0pt] (0,0) -- (1.8,0.5) node[right, lab] {$\mathbf{d}_1$};
  \draw[->, line width=1.0pt] (0,0) -- (0.6,1.7) node[above, lab] {$\mathbf{d}_2$};
  \node[pt] at (1.2,1.1) {};
  \node[lab, anchor=east] at (1.15,1.18) {$\mathbf{x}$};
  \draw[densely dotted] (1.2,1.1) -- (1.8,1.65);
  \node[pt] at (1.8,1.65) {};
  \node[lab, anchor=south east] at (1.85,1.7) {$1.5\,\mathbf{x}$};
  \node[panelcap] at (0,-2.45) {\textbf{cone generated by} $\mathbf{d}_1, \mathbf{d}_2$\\[-1pt]
        $\{t_1\mathbf{d}_1 + t_2\mathbf{d}_2 : t_1, t_2 \ge 0\}$\\[-1pt] a convex cone};
\end{scope}
% ------------------------------------------------------------- right
\begin{scope}[xshift=6.3cm]
  \begin{scope}
    \clip (-2.3,-2.3) rectangle (2.3,2.3);
    \path[pattern=north east lines, pattern color=okabeblue!55]
       (0,0) -- (3,0) -- (3,3) -- (0,3) -- cycle;
    \path[pattern=north east lines, pattern color=okabeblue!55]
       (0,0) -- (-3,0) -- (-3,-3) -- (0,-3) -- cycle;
  \end{scope}
  \draw[okabeblue, line width=1.2pt] (2.3,0) -- (0,0) -- (0,2.3);
  \draw[okabeblue, line width=1.2pt] (-2.3,0) -- (0,0) -- (0,-2.3);
  \draw[ax] (-2.3,0) -- (2.4,0) node[right, lab] {$x_1$};
  \draw[ax] (0,-2.3) -- (0,2.4) node[above, lab] {$x_2$};
  \node[pt] (p) at (1.6,0.4) {};
  \node[pt] (q) at (-0.4,-1.6) {};
  \node[lab, above right] at (p) {$\mathbf{p}$};
  \node[lab, below left] at (q) {$\mathbf{q}$};
  \draw[okabevermillion, line width=1.2pt] (p) -- (q);
  \draw[okabevermillion, line width=1.2pt] (0.5,-0.5) -- (0.7,-0.7);
  \draw[okabevermillion, line width=1.2pt] (0.5,-0.7) -- (0.7,-0.5);
  \node[lab, text=okabevermillion, anchor=west, fill=white, inner sep=1pt] at (0.8,-0.75) {not in $\mathcal{C}$};
  \node[panelcap] at (0,-2.45) {\textbf{a cone that is not convex}\\[-1pt]
        $\{x_1 x_2 \ge 0\}$: the segment from $\mathbf{p}$\\[-1pt] to $\mathbf{q}$ leaves it};
\end{scope}
\end{tikzpicture}
